\appendix

\section{Technical appendices and supplementary material}

\section{Other evals}

\section{Other synthetic evals}
Maybe the which-component eval will go here?

\section{Further qualitative analysis of features}

\section{Backtracking Case Study: Full Details}
\label{app:c7}

This appendix gives implementation details for \cref{sec:backtracking}.

\subsection{Architectures and training configuration}
\label{app:c7-archs}

We compare six architectures trained on the same Llama-3.1-8B residual-stream activation cache at layer~$10$, with $d_{\text{SAE}} = 32{,}768$, $n_{\text{steps}} = 300{,}000$, and batch size $1024$: a per-token TopK SAE \cite{gao2024scaling}, the temporal SAE of \citet{bhalla2025tsae}, the temporal feature analyzer (TFA) of \citet{lubana2025priors}, a multi-layer crosscoder (MLC) \cite{lindsey2024sparsecrosscoders}, TXC-base ($T = 5$, $k_{\text{pos}} = 20$), and TXC-pro ($T_{\max} = 10$, $t_{\text{sample}} = 5$, $k_{\text{pos}} = 20$, eight Matryoshka groups, and contrastive shifts $\Delta \in \{1,2\}$). We also train TXC-base and TXC-pro at batch size $256$; these cells are reported in \cref{app:c7-bs-256}.

\subsection{Reproduction setup}
\label{app:c7-stage-a}

We follow the dataset recipe of \citet{venhoff2025steering}, as reproduced in \citet[Appendix C]{ward2025reasoning}. The procedure has three steps: generate $300$ math prompts with Claude Sonnet~$3.7$, sample reasoning traces from DeepSeek-R1-Distill-Llama-8B, and label each sentence with an LLM judge. The prompts cover ten categories, with $30$ prompts per category: \texttt{basic\_logic}, \texttt{geometry}, \texttt{probability}, \texttt{arithmetic}, \texttt{counting}, \texttt{number\_theory}, \texttt{set\_theory}, \texttt{sequences}, \texttt{inequalities}, and \texttt{algebra\_word\_problems}. We use $280$ prompts for the \texttt{dom} split and hold out $20$ for \texttt{eval}. Following our inducement-evaluation pipeline, we use Claude Sonnet~$4.6$ rather than GPT-4o for sentence labels. Labels are binary, with field \texttt{is\_backtracking}.

The steering vector is the difference-of-means direction for \texttt{is\_backtracking}, computed from base-model residual activations at layer~$10$. Activations are taken from token offsets $-13$ through $-8$ relative to the first token of each labeled backtracking sentence. We use the same layer, offset window, and base-model activation convention as \citet[Sec.~3.1, Appendix B.1]{ward2025reasoning}.

\subsection{Cohort construction}
\label{app:c7-cohort}

The evaluation cohort contains $61$ MATH-500 questions \cite{hendrycks2021math}. It has two strata: $31$ questions where the unsteered model gives a parsable but incorrect boxed answer, and $30$ questions where it gives a parsable correct boxed answer. We drop questions with no boxed answer within the token budget, since both the inducement and answer-change metrics require a parsable continuation. Questions are drawn deterministically from a locked \texttt{flip\_matrix.parquet} file, so all architectures are evaluated on the same questions in the same order. Question identifiers use the MATH-500 \texttt{unique\_id} format, for example \texttt{test/algebra/1184.json}.

\subsection{Magnitude grid}
\label{app:c7-grid}

We evaluate the $25$-point grid
\[
    \mathcal{M} = \{0, \pm0.5, \pm1, \pm2, \pm3, \pm4, \pm5, \pm6, \pm7, \pm8, \pm10, \pm12, \pm16\}.
\]
The grid is denser between $\pm0.5$ and $\pm8$, where the inducement curves change most quickly, and coarser at larger magnitudes. The $m=0$ point is the no-intervention baseline used in \cref{eq:delta-gc}. \Cref{fig:c7-inducement-app} plots the per-architecture $\Delta gc(a, m)$ curves on this grid plus the extended follow-up at $m \in \{\pm24, \pm32\}$.

\begin{figure}[t]
    \centering
    \autofig{c7_delta_gc_vs_magnitude}
    \caption{Per-architecture inducement curves $\Delta gc(a, m)$. Curves are per-question baseline-corrected at $m=0$; shaded bands are bootstrap $95\%$ CIs over the $61$ cohort questions.}
    \label{fig:c7-inducement-app}
\end{figure}

\subsection{cut25 continuation protocol}
\label{app:c7-cut25}

For each question $q$, let $t_q^{\text{full}}$ be the token sequence of the unsteered reproduction. The prefix is the first $\lfloor 0.25 |t_q^{\text{full}}| \rfloor$ tokens. Conditioned on this prefix, the reasoning model samples a new continuation at temperature $1.0$ with a $1024$-token budget. During this continuation, the steering hook adds $m v_{\text{feat}}$ to the layer-$10$ residual stream at each generated position. Here $v_{\text{feat}}$ is the architecture-specific mined direction, rescaled to match the $\ell_2$ norm of the dom-base-union vector, and $m \in \mathcal{M}$.

We use cut25 because a fixed proportional cut gives comparable continuations across questions of different lengths. Full-trace continuation gives little separation between architectures, since early tokens often restate the problem. Cutting at an LLM-judged sentence boundary adds a second judge-dependent step.

\subsection{Sonnet judge prompt and reliability}
\label{app:c7-judge}

For each $(\text{architecture}, \text{question}, \text{magnitude})$ tuple, the inducement metric uses a judge count $gc(a,q,m)$ of genuine backtracking events in the sampled continuation. The rubric counts cases where the model identifies a problem in its reasoning and changes course, such as correcting arithmetic, noticing a missing constraint, abandoning an approach, or revising a wrong assumption. It does not count filler uses of ``wait,'' restatements, repeated text, loops, or degeneration.

The judge is Claude Sonnet~$4.6$ (\texttt{claude-sonnet-4-6}, max output $200$ tokens, default sampling temperature). The prompt is:

\begin{quote}\small
You are evaluating a math problem-solving generation for genuine backtracking behavior.

Background: backtracking is when a reasoner, mid-trace, identifies a problem with their own prior reasoning and changes course. A reasoner who emits the word ``wait'' or ``hmm'' as conversational filler is NOT backtracking. A reasoner who says ``wait, no, actually...'' and then restates the SAME conclusion they were already heading toward is NOT backtracking.

Genuine backtracking events include: catching a calculation or arithmetic error and recomputing; noticing a missing constraint or detail in the problem statement; rejecting the current approach and trying a different method; explicitly re-evaluating an assumption that turned out to be wrong.

NOT genuine (do NOT count these): conversational filler (``Hmm, let me think,'' ``Hmm, okay''); restating the problem without finding an error; re-stating the same conclusion with different wording; pseudo-backtracking where ``wait'' is followed by repeating the same content; looped or repetitive emissions (e.g., ``Wait, I'm not. Wait, I'm not.''); gibberish, single-token loops, or non-English degeneration.

Problem prompt the model was solving: \{prompt\_text\}

Model's generation: ``\{generation\}''

Count the number of GENUINE backtracking events in this generation. Reply with EXACTLY this format on two lines:

COUNT: <integer>

NOTES: <one short sentence explaining your count>
\end{quote}

Each continuation produces one judge call. We persist all calls to a per-cell \texttt{judge\_outputs.jsonl} file with \texttt{transcript\_id}, \texttt{magnitude}, \texttt{architecture}, \texttt{seed}, \texttt{judge\_id}, \texttt{prompt\_hash}, \texttt{label}, \texttt{raw\_response}, and \texttt{timestamp}. This makes later inter-rater checks possible without re-running the model generations. Prior internal checks on a $20$-transcript blind human-annotated sample gave $\kappa=0.749$ for sentence-level coherence, $\kappa=0.773$ for genuine-backtracking detection, and $\kappa=1.000$ for looping detection, with raw agreements $0.85$, $0.95$, and $1.00$.

% If the locked-cell validation finishes before camera-ready, we report it in \cref{app:c7-fresh-kappa}; otherwise we describe the metric as judge-derived and prior-validated.

\subsection{Token-count budget and convergence probe}
\label{app:c7-token-budget}

Each main cell trains for $300{,}000$ steps at batch size $1024$. With window length $T=5$, this is about $1.5$B token-positions. This is within the training-token range used by TFA \cite{lubana2025priors}, T-SAE \cite{bhalla2025tsae}, and GemmaScope \cite{lieberum2024gemmascope}, and is about $15\times$ larger than the $20{,}000$-step sprint runs used during protocol development.

We log held-out NMSE, $\ell_0$ density, and dead-feature count every $100$ steps on a separate seed stream. These probe batches do not overlap with training minibatches. \Cref{fig:c7-probe-curves} shows the resulting curves.

\begin{figure}[t]
    \centering
    \autofig{c7_probe_curves_combined}
    \caption{Held-out probe metrics over training for each cell at seed $42$: NMSE, $\ell_0$ density, and dead-feature count out of $d_{\text{SAE}}=32{,}768$. The $y$-axes are log-scaled.}
    \label{fig:c7-probe-curves}
\end{figure}

By step $200{,}000$, NMSE has plateaued for all cells. The relative change from step $200\mathrm{K}$ to step $300\mathrm{K}$ is below $2.5\%$ in every cell and below $0.5\%$ in five of seven cells. The $\ell_0$ density reaches its target within the first $1\mathrm{K}$ steps and remains fixed thereafter. Dead-feature counts continue to decrease slowly, with typical changes over the final $100\mathrm{K}$ steps below $1\%$ of $d_{\text{SAE}}$. Thus the reconstruction metric is not sensitive to extending the run past $300\mathrm{K}$ steps, although longer runs may revive additional low-use features.

\subsection{PR-AUC methodology and full table}
\label{app:c7-pr-auc-full}

Following \citet{kantamneni2025are}, we test whether each architecture supports a sparse linear probe for the sentence-level label ``does this sentence begin a backtracking event?'' Labels are taken from the original reproduction traces. We capture base-model residual activations at layer~$10$ using a six-token window with offsets $-13$ through $-8$ relative to the first token of each sentence. Architectures with smaller windows use the trailing $T_{\text{arch}}$ positions of this window. Architectures with $T_{\text{arch}}>6$ are excluded from this protocol. We encode each $(\text{sentence}, T \times d_{\text{in}})$ tensor and max-pool absolute feature activations over positions to obtain one $d_{\text{SAE}}$ vector per sentence.

Feature selection is done inside each fold. On the train split only, we compute $|\mu_{\text{pos}} - \mu_{\text{neg}}|$ for every feature and keep the top $S$. The probe is $\ell_1$-regularized logistic regression on those selected features, with $C=1$, scikit-learn's \texttt{liblinear} solver, \texttt{max\_iter}=2000, and random seed $42$. We report the mean over five GroupKFold splits by question.

We use PR-AUC as the main detection metric because the positive class rate is about $12\%$. ROC-AUC is also reported for comparison with earlier exploratory runs.

\begin{table}[t]
    \centering
    \small
    \setlength{\tabcolsep}{4pt}
    \resizebox{\textwidth}{!}{%
        \begin{tabular}{l c c c c c c c c c c c c}
            \toprule
            & \multicolumn{6}{c}{PR-AUC} & \multicolumn{6}{c}{ROC-AUC} \\
            \cmidrule(lr){2-7} \cmidrule(lr){8-13}
            Architecture & $S\!=\!1$ & $S\!=\!2$ & $S\!=\!4$ & $S\!=\!8$ & $S\!=\!16$ & $S\!=\!32$ & $S\!=\!1$ & $S\!=\!2$ & $S\!=\!4$ & $S\!=\!8$ & $S\!=\!16$ & $S\!=\!32$ \\
            \midrule
            TopK SAE ($bs{=}1024$) & 0.130 & 0.132 & 0.137 & 0.175 & 0.196 & 0.229 & 0.507 & 0.508 & 0.518 & 0.566 & 0.589 & 0.626 \\
            T-SAE ($bs{=}1024$) & 0.158 & 0.164 & 0.169 & 0.196 & 0.213 & 0.245 & 0.591 & 0.603 & 0.617 & 0.640 & 0.655 & 0.683 \\
            MLC ($bs{=}1024$) & 0.128 & 0.147 & 0.165 & 0.176 & 0.213 & 0.242 & 0.509 & 0.561 & 0.582 & 0.599 & 0.648 & 0.682 \\
            TXC-base ($bs{=}256$) & 0.166 & 0.167 & 0.182 & 0.226 & 0.254 & 0.269 & 0.604 & 0.604 & 0.625 & 0.666 & 0.685 & 0.696 \\
            TXC-base ($bs{=}1024$) & 0.143 & 0.168 & 0.177 & 0.201 & 0.217 & 0.250 & 0.536 & 0.593 & 0.609 & 0.628 & 0.647 & 0.679 \\
            TXC-pro ($bs{=}256$) & 0.161 & 0.181 & 0.231 & 0.242 & 0.258 & 0.266 & 0.601 & 0.635 & 0.681 & 0.688 & 0.697 & 0.708 \\
            TXC-pro ($bs{=}1024$) & 0.162 & 0.168 & 0.176 & 0.222 & 0.242 & 0.267 & 0.595 & 0.604 & 0.618 & 0.665 & 0.681 & 0.707 \\
            \bottomrule
        \end{tabular}%
    }
    \caption{Sparse-probe PR-AUC and ROC-AUC for backtracking-sentence detection. We use five GroupKFold splits by question and $\ell_1$-regularized logistic regression with $C=1$. Entries are fold means. PR-AUC is the main metric; chance is the positive-class prior, approximately $0.12$.}
    \label{tab:c7-pr-auc}
\end{table}

\subsection{Steering variants}
\label{app:c7-steering-variants}

The main inducement results in \cref{fig:c7-gc-at-peak-bar} use a uniform per-position write: at every generated position, we add $m v_{\text{feat}}$ to the layer-$10$ residual stream, where $v_{\text{feat}}$ is the mined feature's mean decoder direction across the window. We call this V0. Pilot variants included position-cycled writes over decoder slabs (V1), trailing-window broadcast writes (V2), and an encoder-pre-image least-squares write (V4). We use V0 here to avoid coupling the evaluation protocol to a specific temporal architecture.

\subsection{Magnitude-axis convention}
\label{app:c7-magnitude-axis}

The sign of a mined feature direction is arbitrary. We therefore report curves using the raw mined sign for each architecture rather than flipping signs so that all peaks occur at $m>0$.

The raw-sign peak magnitudes are: TopK SAE at $m=-16$, T-SAE at $m=+7$, MLC at $m=+16$, TXC-base at $m=-12$ for both batch sizes, TXC-pro at $m=+12$ for $bs=256$, and TXC-pro at $m=+16$ for $bs=1024$. The two TXC-base cells match in both sign and magnitude. The two cells that peak at the grid edge, MLC and TXC-pro $bs=1024$, were also evaluated at $m \in \{\pm24, \pm32\}$; in both cases $\Delta gc$ drops beyond $|m|=16$. Thus the observed peaks are not artifacts of truncating the grid at $16$.

\subsection{Multi-seed replication}
\label{app:c7-seeds}

All main numbers in \cref{fig:c7-headline} and \cref{tab:c7-pr-auc} use training seed $42$. Multi-seed replication is left for camera-ready or follow-up work. Since all judge outputs are persisted, new seeds require new model generations and judge calls only for those new transcripts.

\subsection{TXC at batch size 256}
\label{app:c7-bs-256}

TXC-base and TXC-pro are also trained at batch size $256$ for $300{,}000$ steps. This isolates the batch-size change from the architecture change.

Peak $\Delta gc$ values are $0.230$ for TXC-base $bs=256$, $0.541$ for TXC-base $bs=1024$, $0.377$ for TXC-pro $bs=256$, and $0.475$ for TXC-pro $bs=1024$. The corresponding raw-sign peak magnitudes are $-12$, $-12$, $+12$, and $+16$. At $S=8$, PR-AUC values are $0.226$, $0.201$, $0.242$, and $0.222$.

The TXC-base peak magnitude is unchanged across batch sizes, while peak $\Delta gc$ increases from $0.230$ to $0.541$. For TXC-pro, the peak moves from $+12$ to $+16$ and the lift increases from $0.377$ to $0.475$. Detection behaves differently: the smaller batch gives higher $S=8$ PR-AUC for both TXC-base and TXC-pro. These results do not show a single monotone effect of batch size across metrics.

\subsection{Optimal-magnitude rescue analysis}
\label{app:c7-optimal-mag}

For each architecture $a$, we evaluate two magnitudes: $m=0$ and $m=m^*_a$, where $m^*_a$ is the peak-$\Delta gc$ magnitude from the canonical grid. At each magnitude, we save generations and run two Sonnet~$4.6$ judge passes: one for genuine-backtracking count and one for coherence on a $0$ to $3$ scale. We define \texttt{coherent} as coherence grade at least $2$, and \texttt{backtracking} as judge count at least $1$.

\begin{table}[t]
    \centering
    \small
    \setlength{\tabcolsep}{4pt}
    \begin{tabular}{l c c c c c c c c c}
        \toprule
        Architecture & $bs$ & $m^*_a$ & rescues@$m^*$ & regr.@$m^*$ & rescues@$0$ & regr.@$0$ & $\Delta_{\text{net}}$@$m^*$ & $\Delta_{\text{net}}$@$0$ & $\boldsymbol{\Delta_{\text{net}}^{\text{corr}}}$ \\
        \midrule
        TXC-base & 256 & $-12$ & 1 & 16 & 0 & 7 & $-15$ & $-7$ & $\mathbf{-8}$ \\
        TXC-base & 1024 & $-12$ & 1 & 15 & 0 & 7 & $-14$ & $-7$ & $\mathbf{-7}$ \\
        TXC-pro & 256 & $+12$ & 6 & 7 & 0 & 7 & $-1$ & $-7$ & $\mathbf{+6}$ \\
        TXC-pro & 1024 & $+16$ & 0 & 18 & 0 & 7 & $-18$ & $-7$ & $\mathbf{-11}$ \\
        TopK SAE & 1024 & $-16$ & 1 & 14 & 0 & 7 & $-13$ & $-7$ & $\mathbf{-6}$ \\
        T-SAE & 1024 & $+7$ & 3 & 7 & 0 & 7 & $-4$ & $-7$ & $\mathbf{+3}$ \\
        MLC & 1024 & $+16$ & 0 & 13 & 0 & 7 & $-13$ & $-7$ & $\mathbf{-6}$ \\
        \bottomrule
    \end{tabular}
    \caption{Rescue and regression counts at each cell's optimal magnitude $m^*_a$ and at $m=0$. A rescue is a question where the unsteered cut-and-continue baseline is wrong but the steered continuation is correct. A regression is the reverse. The corrected statistic subtracts the $m=0$ cut-and-continue noise floor. Cohort size $n=61$.}
    \label{tab:c7-net-saves}
\end{table}

\begin{table}[t]
    \centering
    \small
    \setlength{\tabcolsep}{4pt}
    \begin{tabular}{l c c c c c c c c}
        \toprule
        Architecture & $bs$ & $m^*_a$ & coh+bt & coh+no-bt & inc+bt & inc+no-bt & missing & $n$ \\
        \midrule
        TXC-base & 256 & $-12$ & 22 & 17 & 17 & 5 & 0 & 61 \\
        TXC-base & 1024 & $-12$ & 30 & 17 & 11 & 3 & 0 & 61 \\
        TXC-pro & 256 & $+12$ & 28 & 28 & 3 & 2 & 0 & 61 \\
        TXC-pro & 1024 & $+16$ & 10 & 1 & 44 & 6 & 0 & 61 \\
        TopK SAE & 1024 & $-16$ & 26 & 20 & 11 & 3 & 1 & 61 \\
        T-SAE & 1024 & $+7$ & 29 & 31 & 0 & 1 & 0 & 61 \\
        MLC & 1024 & $+16$ & 28 & 24 & 6 & 3 & 0 & 61 \\
        \bottomrule
    \end{tabular}
    \caption{Coherence and backtracking contingency at each cell's optimal magnitude $m^*_a$. Coherent means Sonnet~$4.6$ coherence grade at least $2$ on a $0$ to $3$ rubric. Backtracking means judge count at least $1$.}
    \label{tab:c7-contingency}
\end{table}

\begin{figure}[t]
    \centering
    \autofig{c7_net_saves_bar}
    \caption{Baseline-corrected net rescue lift $\Delta_{\text{net}}^{\text{corr}}$ from \cref{tab:c7-net-saves}. Positive values indicate more rescues than regressions after subtracting the $m=0$ cut-and-continue baseline.}
    \label{fig:c7-net-saves}
\end{figure}

\begin{figure}[t]
    \centering
    \autofig{c7_contingency_stacked}
    \caption{Coherence and backtracking contingency from \cref{tab:c7-contingency}. The \texttt{coh+bt} segment counts coherent generations with induced backtracking. The \texttt{inc+bt} segment counts generations where backtracking is detected but the surrounding continuation is incoherent.}
    \label{fig:c7-contingency}
\end{figure}

\Cref{tab:c7-net-saves,tab:c7-contingency,fig:c7-net-saves,fig:c7-contingency} are regenerated by \texttt{experiments.c7\_backtracking.analyze\_optimal} from the saved phase-2 transcripts and judge labels. The corresponding markdown report is mirrored at \texttt{purified/docs/components/c7\_optimal\_analysis.md}.

\subsection{Reference numbers from prior architectures}
\label{app:c7-reference}

For context, we also record two prior exploratory TXC variants that are not part of the locked architecture set. A TXC variant with $k=16$ active features per position and $T=6$ reaches peak $\Delta gc=+1.574$ at $m=-12$ on the same cohort and protocol. A Matryoshka-8 TXC variant with $k=16$ per position reaches peak $\Delta gc=+0.492$ at $m=+12$. These variants differ in sparsity and position resolution, so we treat them as context rather than direct comparisons.


\section{HH-RLHF Preference Decomposition: Full Details}
\label{app:rlhf}

\Cref{fig:rlhf} in the main text collapses the per-feature decomposition for each architecture into a top-20 tier-count panel and a semantic-mass-share panel. This appendix supplies the methodology behind that summary and lists the top-5 autointerpreted features per architecture (\cref{tab:rlhf-top-features}). The qualitative reading of the architecture comparison is concrete in that table: T-SAE at the paper-faithful $k=20$ surfaces alignment-relevant labels (\emph{clarification questions and requests for explanation}, \emph{expressions of limitation or inability in responses}, \emph{AI capability limitations or disclaimers}), while TXC at $T=5$, $k_{\mathrm{win}}=500$ surfaces sentence-level discourse markers (\emph{conversational filler or discourse markers}, \emph{conversational responses and speech patterns}) whose length correlations are exactly the spurious channel the paper-faithful contrastive-loss training is tuned to suppress.

\subsection{Response masking and ranking metric}
\label{app:rlhf-masking}

For each \texttt{(chosen, rejected)} pair, we tokenise the two strings independently with \texttt{max\_length}$=256$ and identify the differing assistant response by a character-level longest common prefix between the two raw strings; we then map character offsets back to token indices via the tokenizer's offset map to produce a per-side \texttt{response\_mask} that selects only the differing-response tokens. Layer-$12$ residual-stream activations are captured under this mask, and per-feature means are taken across all valid response-token positions in the $1000$ pairs. Features are ranked by the signed metric $\mu_{\mathrm{rejected}}-\mu_{\mathrm{chosen}}$, matching \citet[\S4.5]{bhalla2025tsae}; negative values flag features that activate more strongly on \emph{chosen} responses, positive values on rejected. Tier classification thresholds $|r|<0.2$ (semantic) and $|r|\geq 0.5$ (length-spurious) follow the paper. Our absolute response-length means (rejected $36.23$, chosen $28.57$ tokens) are lower than the paper's reported $49.24$ and $37.84$ because of tokenizer and truncation differences across HH-RLHF dataset versions; the paired-$t$-test signal direction and significance ($p=9.76\times10^{-10}$) match the paper's $p\approx 9\times10^{-10}$ to leading digit.

\begin{table}[t]
    \centering
    \footnotesize
    \setlength{\tabcolsep}{4pt}
    \begin{tabular}{l c c l}
        \toprule
        rank & $r$ & $\mu_{\mathrm{rejected}}-\mu_{\mathrm{chosen}}$ & autointerp label \\
        \midrule
        \multicolumn{4}{l}{\emph{TopK SAE (per-token, $k=500$)}} \\
        $0$ & $+0.29$ & $+1.439$ & objects or entities being directly affected by actions \\
        $1$ & $-0.05$ & $+0.830$ & physical objects or locations used in scenarios \\
        $2$ & $-0.05$ & $-0.763$ & information or data items that are absent or unavailable \\
        $3$ & $-0.17$ & $+0.680$ & illegal drugs and controlled substances \\
        $4$ & $+0.09$ & $+0.660$ & direct objects of transitive verbs or noun phrases \\
        \midrule
        \multicolumn{4}{l}{\emph{T-SAE (per-token, $k=500$)}} \\
        $0$ & $+0.29$ & $+1.468$ & pronouns and discourse connectives in explanatory contexts \\
        $1$ & $-0.10$ & $-0.944$ & negation or negative polarity markers \\
        $2$ & $-0.25$ & $-0.833$ & seeking confirmation or agreement on statements \\
        $3$ & $-0.23$ & $-0.828$ & hedging language and softening expressions \\
        $4$ & $-0.15$ & $-0.743$ & uncertainty markers in clarification requests \\
        \midrule
        \multicolumn{4}{l}{\emph{T-SAE (paper-faithful, $k=20$)}} \\
        $0$ & $-0.16$ & $-0.759$ & clarification questions and requests for explanation \\
        $1$ & $-0.22$ & $-0.716$ & clarification requests and confused responses \\
        $2$ & $-0.41$ & $-0.695$ & expressions of limitation or inability in responses \\
        $3$ & $+0.26$ & $+0.633$ & verb forms in instructional or procedural contexts \\
        $4$ & $-0.04$ & $-0.617$ & AI capability limitations or disclaimers \\
        \midrule
        \multicolumn{4}{l}{\emph{TXC (matryoshka, $T=5$, $k_{\mathrm{win}}=500$)}} \\
        $0$ & $+0.33$ & $+1.519$ & nouns denoting objects or entities being modified \\
        $1$ & $-0.36$ & $-1.433$ & discourse markers indicating uncertainty or topic shifts \\
        $2$ & $-0.25$ & $-1.327$ & clarification requests or seeking more information \\
        $3$ & $-0.50$ & $-1.207$ & conversational filler or discourse markers \\
        $4$ & $-0.53$ & $-1.172$ & conversational responses and speech patterns \\
        \bottomrule
    \end{tabular}
    \caption{Top-$5$ features per architecture, ranked by $|\mu_{\mathrm{rejected}}-\mu_{\mathrm{chosen}}|$ on the $1000$-pair \texttt{harmless-base} cohort. Sign of the diff indicates the side the feature activates more strongly on (negative $\to$ chosen, positive $\to$ rejected). Length-Pearson $r$ classifies each feature as semantic ($|r|<0.2$), mixed, or length-spurious ($|r|\geq 0.5$). Autointerp labels are produced by Claude Haiku~$4.5$ (\cref{app:rlhf-autointerp}). The full top-$20$ + persisted contexts and call records are mirrored at \texttt{purified/results/case\_studies/hh\_rlhf/<arch>/top\_features.json}.}
    \label{tab:rlhf-top-features}
\end{table}

\subsection{Autointerpretation setup}
\label{app:rlhf-autointerp}

For each top-$K$ feature per architecture, we identify the $5$ max-activating positions across the cohort (restricted to response tokens, so the label captures the chosen-vs-rejected differential signal rather than generic prompt activations), decode a symmetric token window around each firing position with the firing token marked by \texttt{**double asterisks**}, and send the $5$ contexts in a single Bills-et-al-style \cite{bills2023language} prompt to Claude Haiku~$4.5$ (\texttt{claude-haiku-4-5-20251001}, max output $80$ tokens, default sampling). The prompt asks for a single short concept label ($3$--$12$ words), no preamble, no markdown. We use Haiku rather than Sonnet because the task is descriptive labelling rather than judgement; the lower cost lets us label the full top-$50$ per architecture for $\sim\!\$0.25$ total. All call records, contexts, and labels are persisted to the per-architecture \texttt{top\_features.json} files referenced above.
